\begingroup
\makeatletter\tagsleft@true\let\veqno\@@leqno\makeatother
\ArithmeticHead{274}{L}{ANALYSIS. NACHLASS.}
\SourceLeaf{281b}
\pdfbookmark[1]{Erläuterungen zu [V.]}{x1-agm-erlaeuterung-v}
\begin{center}\textls[120]{Erläuterungen zu [V.], S. 194—206.}\end{center}
\par\indent Die Funktion [2] ist der in der Tagebuchnotiz Nr. 108 von Ende Mai 1800 erwähnte »sinus lemniscaticus universalissime acceptus«\textsuperscript{**)}. Sie ist mit den in [1] durch das agM. definierten Größen $\varpi,\varpi'$ nach der Analogie des sinus lemniscaticus gebildet. In der Tat haben wir z. B. in der Scheda Aa, S. 8 die Darstellung
\begin{equation*}\tag*{$[2]'$}\operatorname{sl}\varphi=\frac\pi\varpi\frac4{e^{\frac12\pi}+e^{-\frac12\pi}}\operatorname{sc}\varphi-\frac\pi\varpi\frac4{e^{\frac32\pi}-e^{-\frac32\pi}}\operatorname{sc}3\varphi\ \text{etc},\end{equation*}
\noindent wo sl den sinus lemniscaticus, sc den sinus circuli, d. h. den gewöhnlichen Sinus bedeutet und $\varphi$ beiderseits im \textls[120]{Gradmaß} zu denken ist\textsuperscript{***)}; will man also, wie es \textsc{Schering} beim Abdruck dieser und ähnlicher Formeln im Bande III der Werke (S. 417 ff.) getan hat, auf \textls[120]{Bogenmaß} reduzieren, so hat man für $\varphi=180^0$ linker Hand den halben Lemniskatenumfang $\varpi$ und rechter Hand den halben Kreisumfang, also $\pi$ zu nehmen, dabei ist das lemniskatische $\varpi$ aus dem in [1] definierten $\varpi$ zu bilden, indem man $\mu=1$ setzt. Es ist nämlich (siehe die Tagebuchnotiz Nr. 98, oben S. 260) in diesem Fall
\begin{equation*}\tag*{$[1]'$}\varpi=\varpi'=\frac\pi{M(\surd2,1)}.\end{equation*}
\noindent Für diesen Wert $\mu=1$ ergibt die Formel [2] unmittelbar die Darstellung $[2]'$ des sinus lemniscaticus in der Form, wie sie Werke III, S. 417, art. [6.] angegeben ist\textsuperscript{\(\dagger\))}, d. h. beiderseits auf Bogenmaß reduziert.
\Needspace{4\baselineskip}
\par\smallskip\noindent\rule{.25\linewidth}{.4pt}\par
{\small
\par\indent **) An einer späteren Stelle der Scheda Ac, nämlich S. 39, siehe den art. [8] oben S. 200, bezeichnet \textsc{Gauss} die Umkehrungsfunktion des Integrals
\[\varphi=\int\frac{dx}{\sqrt{(1-x^2)\left(1+\frac{x^2}{\mu^2}\right)}}\]
\noindent geradezu mit $x=\operatorname{sl}\varphi$, also mit demselben Zeichen, das er auch häufig (siehe z. B. die gleich folgende Gleichung $[2]'$ in unserem Text) für den sinus lemniscaticus im engern Sinne verwendet.
\par\indent ***) Vergl. die Aufzeichnung oben S. 170, art. [8.], wo aber $\varphi$ und $\chi$ im Bogenmaß zu denken sind, was durch den Faktor $\frac\pi\varpi$ bei $\varphi$ angezeigt wird.
\par\indent $\dagger$) \textsc{Schering} hat in dieser Formel $\psi$ an Stelle des \textsc{Gauss}schen $n$ gesetzt; mit derselben Änderung hat er auch die Formel [2] und einen Teil der übrigen in unserem [V.] Abschnitt zusammengestellten Formeln Werke III, S. 433—435 wiedergegeben.
\par}
\clearpage
\fancyhead{}\fancyhead[R]{\small275}\fancyhead[C]{\scriptsize BEMERKUNGEN ZUM AGM. ABSCHNITT V.}
\SourceLeaf{282}
\par\indent Im art. [2.] werden die Funktionen $T$ und $W$ eingeführt; sie sind ganz nach der Analogie der lemniskatischen $P,Q$ gebildet, die \textsc{Gauss} in der Scheda Aa eingeführt hatte (siehe Werke III, S. 415 art. [4.] ff., insbesondere art. [8.], S. 418) und die auch in der Scheda Ac (siehe in den artt. [8.], S. 201 und [11.], S. 203 die Formeln für $P^2$ und $Q^2$) vorkommen. Hier bei den Funktionen $T$ und $W$ tritt auch wieder $\varphi$ im Gradmaß auf, was die Gleichung $T\,90^0=\surd\cos\nu$ besonders deutlich zeigt. Es ist jedoch zu bemerken, daß \textsc{Gauss} bei dieser Bezeichnungsweise links, d. h. bei $T\varphi,W\varphi$, geradezu $\varphi=n\varpi$, und rechts, d. h. unter dem trigonometrischen Funktionszeichen ebenso $\varphi=n\pi$ nimmt, also nicht, wie wir erwarten würden, $\varphi^0=n\,.\,180^0$.
\par\indent Um die Beziehung zwischen den Funktionen $T,W$ und den \textsc{Jacobi}schen Thetafunktionen möglichst deutlich hervortreten zu lassen, sollen die in den artt. [2.], [3.] und [5.] enthaltenen \textsc{Gauss}schen Formeln hier in der neueren Bezeichnungsweise wiedergegeben werden, wobei wir uns an H. \textsc{Webers} \textit{Elliptische Funktionen und algebraische Zahlen}, Braunschweig 1891, anschließen wollen.
\par\indent Setzt man
\[q=e^{-\frac{\varpi'}\varpi\pi}\]
\noindent und
\[\begin{aligned}
\vartheta_{00}\,n|q&=1+2q\cos2n\pi+2q^4\cos4n\pi+2q^9\cos6n\pi+\ldots\\[4pt]
\vartheta_{11}\,n\,q)&=2q^{\frac14}\sin n\pi-2q^{\frac94}\sin3n\pi+2q^{\frac{25}4}\sin5n\pi-\ldots\\[4pt]
\vartheta_{10}(n|q)&=2q^{\frac14}\cos n\pi+2q^{\frac94}\cos3n\pi+2q^{\frac{25}4}\cos5n\pi+\ldots\\[4pt]
\vartheta_{01}(n|q)&=1\quad2q\cos2n\pi+2q^4\cos4n\pi-2q^9\cos6n\pi+\ldots
\end{aligned}\]
\noindent so ist, wenn wir in Übereinstimmung mit \textsc{Gauss} $\varphi=n\varpi$ nehmen,
\[\begin{aligned}
W\varphi&=\sqrt{M(1,\cos\nu)}\,\vartheta_{00}(n|q)\\[4pt]
T\varphi&=\frac1{\sqrt\mu}\sqrt{M(1,\cos\nu)}\,\vartheta_{11}(n|q)\\[4pt]
T(\tfrac12\varpi-\varphi)&=\frac1{\sqrt\mu}\sqrt{M(1,\cos\nu)}\,\vartheta_{10}(n|q)\\[4pt]
W(\tfrac12\varpi-\varphi)&=\sqrt{M(1,\cos\nu)}\,\vartheta_{01}(n|q\,.
\end{aligned}\]
\noindent Die Formeln [10], [12] geben also $\vartheta_{11}^2(n|q),\vartheta_{00}^2(n|q)$ dargestellt durch $\vartheta_{00}(2nq^2)$ und $\vartheta_{10}(2n|q^2)$, die Formeln [14], [15] sind den folgenden (siehe \textsc{Weber} a. a. O., S. 79, 80) äquivalent:
\begin{equation*}\tag*{$[14]'$}\left\{\begin{aligned}
2\vartheta_{00}(0|q^2)\vartheta_{00}(2n|q^2)&=\vartheta_{00}^2(n|q)+\vartheta_{01}^2(n|q),\\[-1pt]
&\hspace{1em}\text{(\textsc{Weber} (16), S. 80)}\\[3pt]
2\vartheta_{00}(0|q^2)\vartheta_{10}(2n|q^2)&=\vartheta_{10}^2(n|q)-\vartheta_{11}^2(n|q),\\[-1pt]
&\hspace{1em}\text{(\textsc{Weber} (14), S. 80)}
\end{aligned}\right.\end{equation*}
\noindent wobei die Gleichungen [16], die die Quadrate der Funktionen $\vartheta_{10}(n|q),\vartheta_{01}(n|q)$ durch die der Funktionen $\vartheta_{00}(n|q),\vartheta_{11}(n|q)$ darstellen (vergl. \textsc{Weber} a. a. O., S. 54), heranzuziehen sind. Die Gleichungen $[14]'$ stellen also für die Thetafunktionen die sogenannte \textsc{Landen}sche Transformation dar. Demgegenüber gibt die Gleichung [17] die Darstellung von $\vartheta_{11}(n|q^{\frac12})$ durch das Produkt $\vartheta_{11}(n|q)\,\vartheta_{01}(n|q)$, nämlich:
\begin{equation*}\tag*{$[17]'$}\begin{gathered}
\vartheta_{10}(0|\sqrt q)\,\vartheta_{11}(n|\sqrt q)=2\vartheta_{01}(n|q)\,\vartheta_{11}(n|q),\\
\text{(\textsc{Weber} (10), S. 79)}
\end{gathered}\end{equation*}
\noindent also die sogenannte \textsc{Gauss}sche Transformation (vergl. die Gleichung $[6]'$ unten S. 286.).
\par\smallskip\hfill35\textsuperscript{*}\par
\clearpage
\ArithmeticHead{276}{L}{ANALYSIS. NACHLASS.}
\SourceLeaf{283}
\par\indent In Bezug auf die Konstanten haben wir
\begin{equation*}\tag*{$(\alpha)$}\begin{gathered}
\frac1{M(1,\cos\nu)}=\vartheta_{00}^2(0|q),\qquad\frac{\cos\nu}{M(1,\cos\nu)}=\vartheta_{01}^2(0|q),\\
\frac{\sin\nu}{M(1,\cos\nu)}=\vartheta_{10}^2(0|q);
\end{gathered}\end{equation*}
\noindent die dritte dieser Gleichungen ist nichts anderes als die zweite Gleichung von [3]
\[T\,90^0=T\frac\varpi2=\sqrt{\cos\nu},\]
\noindent die im wesentlichen auf die im \textsc{Leiste} aufgezeichnete Gleichung [6] des Abschnitts [II.], oben S. 177, hinauskommt, wenn man beachtet, daß unser $q$ gleich dem dortigen $z^2$ ist. —
\par\indent Ferner gelten die Relationen
\begin{equation*}\tag*{$(\beta)$}\left\{\begin{aligned}
\vartheta_{00}^2(0|q^2)+\vartheta_{10}^2(0|q^2)&=\vartheta_{00}^2(0|q)\\
\vartheta_{00}^2(0|q^2)-\vartheta_{10}^2(0|q^2)&=\vartheta_{01}^2(0|q)\\
2\vartheta_{00}(0|q^2)\,\vartheta_{10}(0|q^2)&=\vartheta_{10}^2(0|q),
\end{aligned}\right.\end{equation*}
\noindent die für
\[\begin{array}{l|l}
a_0=M(a_0,b_0)\,\vartheta_{00}^2(0|q)&a_1=M(a_0,b_0)\,\vartheta_{00}^2(0|q^2)\\
b_0=M(a_0,b_0)\,\vartheta_{01}^2(0|q)&b_1=M(a_0,b_0)\,\vartheta_{01}^2(0|q^2)\\
c_0=M(a_0,b_0)\,\vartheta_{10}^2(0|q)&c_1=M(a_0,b_0)\,\vartheta_{10}^2(0|q^2)
\end{array}\]
\noindent den Relationen des agM.:
\[a_0=a_1+c_1,\quad b_0=a_1-c_1,\quad c_0=2\sqrt{a_1c_1}\]
\noindent entsprechen. Mit den in unserem \textit{Abriß}, Artikel 3., S. 255 benutzten \textsc{Gauss}schen Bezeichnungen ist
\[\vartheta_{00}(0|x)=p(x),\quad\vartheta_{01}(0|x)=q(x),\quad\vartheta_{10}(0|x)=r(x).\]
\noindent Die Bezeichnung der Gln. [19], [20] des art. [5.] weicht von dieser etwas ab.
\par\indent Charakteristisch für die hier umrissene Theorie der allgemeinen elliptischen Funktionen ist, daß \textsc{Gauss} nicht den in dem \textit{Problema} zu Beginn des art. [1.], S. 194 formulierten Gedankengang weiter verfolgt, sondern, von dem elliptischen Integral völlig losgelöst, nach der Analogie des sinus lemniscaticus die Funktion $S\varphi$, den sinus lemniscaticus universalissime acceptus, mit den beiden Perioden $2\varpi$ und $2\varpi'i$ bildet, die er wiederum durch die Gleichungen [1] mit Hilfe des agM. definiert hat. In der oben S. 248 abgedruckten Briefstelle an \textsc{Bessel} (vom 30. März 1828) sagt \textsc{Gauss}, \textsc{Abel} habe denselben Weg genommen, den er (\textsc{Gauss}) 1798 eingeschlagen hatte; gemeint sind hier offenbar die auf lemniskatische Funktionen bezüglichen Arbeiten aus der zweiten Hälfte des Jahres 1798 (Juli-Oktober), wie sie in der Scheda Aa aufgezeichnet (gedruckt Werke III, S. 413—419, artt. [1.]—[4.], [6.]—[8.]) und in den Tagebuchnotizen Nr. 92 und Nr. 95 angezeigt sind. Zugleich geht aber aus diesem Ausspruche von \textsc{Gauss} hervor, daß er bei der Entwicklung der allgemeinen Theorie der elliptischen Funktionen (denn um diese handelt es sich ja bei \textsc{Abel}) ganz nach der Analogie der lemniskatischen Funktionen vorgegangen ist. In der Tat sehen wir auch in den in diesem Abschnitte wiedergegebenen Aufzeichnungen zwischen den auf die allgemeinen elliptischen Funktionen bezüglichen Formeln immer wieder lemniskatische Formeln auftreten, so im art. [2.] und weiterhin in den artt. [7.], [8.], [11.]. Auch die Schlußbemerkung dieses Abschnitts — mit der auch die Scheda Ac abschließt — zeigt, daß \textsc{Gauss} die ganze hier behandelte Funktionsklasse als in der Betrachtung der Lemniskate wurzelnd nach dieser Kurve benannt hat.
\par\indent Zur Herleitung der im art. [3.] zusammengestellten Formeln für die Transformation zweiter Ordnung der Thetafunktionen, d. h. um die Quadrate der vier Funktionen
\[T\varphi,\ W\varphi,\ W\left(\frac\varpi2-\varphi\right)\ T\left(\frac\varpi2-\varphi\right)\]
\clearpage
\fancyhead{}\fancyhead[R]{\small277}\fancyhead[C]{\scriptsize BEMERKUNGEN ZUM AGM. ABSCHNITT V.}
\SourceLeaf{284}
\noindent und das Produkt $T\varphi\,W\left(\frac\varpi2-\varphi\right)$ darzustellen, scheint \textsc{Gauss} sich schon hier eines Prinzips bedient zu haben, das man jetzt als das Transformationsprinzip von \textsc{Hermite} bezeichnet\textsuperscript{*)}. Mit voller Deutlichkeit tritt dieses Prinzip freilich erst im art. [12.] des weiter unten aus dem Handbuch 16, Bb abgedruckten Abschnitts [I.] der Abhandlung »\textit{Zur Theorie der transscendenten Functionen gehörig}« hervor, vergl. auch den art. [5.] dieses Abschnitts, wo \textsc{Gauss} die Transformationsformeln zusammenstellt und darüber schreibt »aus einem allgemeinen Theorem abzuleiten, s. u.«, womit eben jener Satz des art. [12.] gemeint ist. Aus diesen Transformationsformeln ergeben sich für die Nullwerte
\[W0,\ W\tfrac12\varpi,\ T\tfrac12\varpi\]
\noindent Beziehungen, die ohne weiteres den Algorithmus des agM. liefern, also zu einer Bestätigung der Formeln dienen, die \textsc{Gauss} im art. [5.], oben S. 197, zusammengestellt hat\textsuperscript{**)}. Dieser Artikel enthält dem art. [7.] des Abschnitts [IV.] gegenüber nichts wesentlich Neues; bemerkenswert ist nur die explizite Darstellung der Größe $z$ in den Gln. [18] und [23] und die auch weiterhin im art. [8.], S. 200, noch einmal auftretende Gleichung $\sin\psi=\operatorname{tg}^2\frac\varphi2$, die den Modul der \textsc{Landen}schen Transformation liefert.
\par\indent Der art. [4.], S. 196, gibt das Additionstheorem für das Integral
\[\varphi=\int\frac{dS}{\sqrt{(1-S^2)(1+\mu^2S^2)}};\]
\noindent durch die Transformation $S=\sin u$ verwandelt sich dieses Integral in
\[\int\frac{du}{\sqrt{1+\mu^2\sin^2u}},\]
\noindent das \textsc{Gauss} auf das Titelblatt der Scheda Ac gesetzt hat, siehe S. 184 in der Überschrift.
\par\indent Die lemniskatischen Formeln im art. [7.] (die übrigens ebenso wie die der artt. [2.], [8.] und [11.] zum Teil bereits Werke III, S. 419, 420 abgedruckt sind) bedürfen keiner Erläuterung. Im art. [8.] bedeutet sl — wie bereits oben bemerkt — den sinus lemniscaticus universalissime acceptus. Die am Schluß von art. [8.] stehenden Formeln für die lemniskatischen $(P\varphi)^2,(Q\varphi)^2$ werden im art. [11.] zum Teil wiederholt und numerisch ausgeführt.
\par\indent Bei der im art. [9.] behandelten Umformung von Potenzreihen, deren Exponenten eine arithmetische Progression zweiter Ordnung bilden, in Produkte konnte \textsc{Gauss} an \textsc{Euler} und an seine eigenen älteren Untersuchungen anknüpfen. Unter diesen letzteren erwähnen wir hier nur den art. [9.] der \textit{Exercitationes mathematicae}, oben S. 142, und die S. 26 der Scheda Aa (also im Spätsommer 1798) formulierte Aufgabe:
\par\indent{\large Investigatio factorum progressionis infinitae}\par
\[1+x+x^3+x^6+x^{10}+x^{15}+x^{21}+\text{etc.}=S.\]
\noindent Die Lösung dieser Aufgabe gibt die Gl. 3) des art. [9.]; sie findet sich auch in der 1808 veröffentlichten \textit{Summa io quarundam serierum singularium}, Werke II, S. 20. Das hier mit $\Pi z$ bezeichnete Produkt hat \textsc{Gauss} später (siehe den art. [4.] des Abschnitts [I.] der weiter unten folgenden Abhandlung \textit{Zur Theorie der}
\Needspace{4\baselineskip}
\par\smallskip\noindent\rule{.25\linewidth}{.4pt}\par
{\small
\par\indent *) Siehe \textsc{Hermites} zweiten Brief an \textsc{Jacobi} vom August 1844, \textsc{Crelles} Journal für Mathematik \textbf{32}, 1845, S. 176, Oeuvres de \textsc{Ch. Hermite}, I (1905), S. 18 ff., besonders S. 24, 25.
\par\indent **) Vergl. die Tagebuchaufzeichnung Nr. 106 vom 22. Mai 1800, wo es heißt: »per quod simul omnia praecedentia [d. h. die Theorie der lemniskatischen Funktionen] nec non theoria mediorum arithmetico-geometricorum pulcherrime nectuntur infinitiesque augentur.«
\par}
\clearpage
\ArithmeticHead{278}{L}{ANALYSIS. NACHLASS.}
\SourceLeaf{285a}
\noindent\textit{transscendenten Functionen gehörig} aus Handbuch 16, Bb und die Werke III, S. 446 ff. aus Handbuch 18, Bd, S. 221—233 abgedruckte Abhandlung) mit $[z]$ und (siehe art. 5 der \textit{Hundert Theoreme über die neuen Transscendenten}, Werke III, S. 465) mit $Fz$ bezeichnet. Die Buchstaben $p,q$ bedeuten hier dasselbe wie im art. [5.], für ihre hier gegebenen Ausdrücke vergl. Werke III, S. 446 ff. aus Handbuch 18, Bd.
\par\indent Im art. [10.] treten als durchaus neues Element die von zwei Veränderlichen $a,z$ abhängigen Produkte auf, ihre Umwandlung in eine Reihe erfolgt nach einer Methode, die später auch \textsc{Jacobi} (1829, \textsc{Jacobis} Werke I, S. 232) ersonnen hat. Im art. [12.] wird so das Produkt $K$, das die allgemeine Thetafunktion darstellt, und sein Quadrat in eine Reihe verwandelt, wobei aber die konstanten Faktoren $P$ bezw. $P$ und $Q$ unbestimmt bleiben. Nimmt man bei der Umformung von $K$ noch die Bestimmung von $P=\frac1{\Pi x^2}$ hinzu, so erhält man die \textls[120]{berühmte Identität} zwischen der Reihen- und Produktdarstellung der Thetafunktionen in der Form, wie sie \textsc{Gauss} später (siehe Abschnitt [I.], art. [4] der unten folgenden Abhandlung \textit{Zur Theorie der transscendenten Functionen gehörig}, ferner \textit{Hundert Theoreme}, Werke III, S. 464, das vierte Theorem) wiederholt angegeben hat. Bekanntlich hat \textsc{Jacobi} in einem Briefe an P. H. \textsc{Fuss}\textsuperscript{*)} erklärt, es sei diese Identität »wohl das wichtigste und fruchtbarste«, was er »in reiner Mathematik erfunden habe«. Die artt. [13.], [14.] geben die Anwendungen auf die Funktionen $T$ und $W$. Im art. [14.] bedeutet $\varepsilon$ ebenso wie im art. [8.] die $\sqrt{-1}$. Wir bemerken noch, daß sich zwischen dem art. [13.] und dem die Scheda Ac abschließenden art. [14.] auf S. 47 beginnend eine Anwendung der Theorie der elliptischen Funktionen findet: \textit{Motus solidi a nullis viribus sollicitati}, über die \textsc{Schering} Werke III, S. 495 kurz berichtet.
\par\indent Die hier besprochenen Aufzeichnungen entsprechen den Tagebuchnotizen Nr. 105—Nr. 108, sie sind demnach im Mai 1800 in Braunschweig verfaßt worden. Man vergl. hierzu den Abschnitt V. des Aufsatzes »Über \textsc{Gauss}' Arbeiten zur Funktionentheorie«.
\Needspace{4\baselineskip}
\par\smallskip\noindent\rule{.25\linewidth}{.4pt}\par
{\small
\par\indent *) \textit{Briefwechsel zwischen Jacobi und Fuss}, herausg. von \textsc{Stäckel} und \textsc{Ahrens}, 1908, S. 60, siehe auch Bibliotheca Mathematica, (3) \textbf{8}, S. 291.
\par}
\clearpage
\endgroup

